ASMPT develops semiconductor assembly and packaging equipment, including high-precision motion systems for advanced packaging. Advanced packaging refers to combining multiple semiconductor functions within a single package, for example through heterogeneous integration of different die types such as logic and memory, and through hybrid bonding, in which dies are directly joined at very fine pitch. Such technologies place strict demands on alignment accuracy, yield, throughput, and manufacturing cost \cite{irds2022,workman2021diewafer}.
\\
This project focuses on a dual-gantry motion system in which two parallel linear actuators support a common beam carrying the payload. The shared beam couples the actuator dynamics, so motion of one actuator influences the other, and the system exhibits position-dependent behaviour and flexible dynamics \cite{garcia2013model, teo2007dynamic}. In such systems, settling behaviour is a key performance metric. After each motion, the system must wait for residual oscillations to decay before the next placement can begin \cite{lambrechts2005trajectory}, so settling time directly affects throughput. Furthermore, cross-coupling means that the motion of one gantry can disturb the other, making accurate prediction relevant not only for settling but also for tracking behaviour during simultaneous motion.
\\
Improving settling performance therefore requires an accurate representation of the underlying plant dynamics. ASMPT currently captures part of the position dependence of the dual-gantry system using a baseline model based on frequency-response measurements at multiple operating points with interpolation across the operating space. However, this baseline captures only part of the position dependence and does not systematically account for cross-coupling between the gantries. A more accurate plant model is therefore needed. Such models are not only valuable for settling prediction, but also for digital-twin development, offline evaluation, control design, and condition monitoring. Digital twins can support offline simulation of semiconductor manufacturing equipment and help optimize operating parameters before testing changes on the real machine, reducing costly trial-and-error \cite{hwang2024semiconductor}. Accurate models can also support condition monitoring and early fault detection, as demonstrated in robotic manufacturing systems \cite{aivaliotis2019predictive}.
\\
White-box models are attractive because they reflect the physics of the system, which makes them interpretable. First-principles (FP) models are a specific type of white-box model, derived from physical laws and described by quantities such as mass, stiffness, and damping. However, purely physics-based models do not always capture all behaviour needed for accurate prediction across the operating range \cite{willard2020integrating}. For the dual-gantry system, this limitation is also evident in the first-principles model structure of Garc\'ia-Herreros et al. \cite{garcia2013model}, which is only able to capture part of the position-dependent behaviour and cross-coupling.
\\
Purely data-driven models such as neural networks or Gaussian processes can describe measured behaviour well, but they give less physical insight, are often harder to trust at operating conditions not covered by the training data \cite{willard2020integrating}. Typically these models also require large amounts of data and significant computational power to train. Physics-informed approaches, including physics-informed neural networks (PINNs), enforce physical laws as constraints during training \cite{karniadakis2021physicsinformed, raissi2017physicsinformedpart1}. For the plant model of the dual-gantry system, both predictive accuracy and physical interpretability are required. However, neither a purely physics-based nor a purely data-driven model is sufficient on its own \cite{willard2020integrating, noel2018greybox}. Grey-box modeling is the logical choice, because it is able to combine a physically meaningful baseline model needed for interpretation with data-driven flexibility needed for accurate prediction \cite{noel2018greybox}. However, related work at ASMPT shows that preserving physical interpretability during model updating remains an open challenge \cite{kessels2025ai}.
\\
Learning-based hybrid modeling has already been demonstrated on real industrial systems \cite{retzler2024}. However, a key gap remains: there is no validated grey-box modeling framework for the dual-gantry system that simultaneously addresses coupled multi-input multi-output (MIMO) dynamics, position-dependent behaviour, and preserves physical interpretability. This project therefore aims to develop and validate such a framework so that the plant model becomes more accurate for settling prediction while remaining physically meaningful.
